\begin{titlepage}


% \begin{center}

% \textsc{\LARGE \univname}\\[1.5cm] % University name
% \textsc{\Large Doctoral Thesis}\\[0.5cm] % Thesis type
%
% \HRule \\[0.4cm] % Horizontal line
% {\huge \bfseries \ttitle}\\[0.4cm] % Thesis title
% \HRule \\[1.5cm] % Horizontal line
%
% \begin{minipage}{0.4\textwidth}
% \begin{center} \large
% % \emph{Author:}\\
% \href{}{\authorname} % Author name - remove the \href bracket to remove the link
% \end{center}
% \end{minipage}\\[3cm]
%


% \large \textit{A thesis in fulfillment of the requirements\\ for the degree of \degreename}\\[0.3cm] % University requirement text
% \textit{in the}\\[0.4cm]
% \groupname\\\deptname\\[2cm] % Research group name and department name
%
% {\large \today}\\[4cm] % Date
% %\includegraphics{Logo} % University/department logo - uncomment to place it
%
% \vfill
% \end{center}


\begin{center}
	VRIJE UNIVERSITEIT

	\vspace{1cm}

	Controlling Penguins

	\vspace{2cm}

	ACADEMISCH PROEFSCHRIFT

	\vspace{0.5cm}

\begin{minipage}{0.65\textwidth}
\centering
ter verkrijging van de graad Doctor aan\\
de Vrije Universiteit Amsterdam,\\
op gezag van de rector magnificus\\
prof.dr. V. Subramaniam,\\
in het openbaar te verdedigen\\
ten overstaan van de promotiecommissie\\
van de Faculteit der Exacte Wetenschappen\\
op dinsdag 3 oktober 2017 om 11.45 uur\\
in het auditorium van de universiteit,\\
De Boelelaan 1105
\end{minipage}\\[3cm]

door\\
\vspace{0.5cm}
Vasileios Syropoulos\\
\vspace{0.5cm}
geboren te Thessaloniki, Griekenland

\end{center}

\clearpage

\begin{flushleft}
promotor: prof.dr. H.G. Raven\\
copromotor: dr. P. Koppenburg
\end{flushleft}

\vfill


\clearpage

\begin{center}
	Cover designed by Antwnis Karanautis\\

	\vspace{1cm}

%	Typeset with \LaTeX \\
	\vspace{0.1cm}
	Printed by Gildeprint - The Netherlands\\
	ISBN: 978-94-6233-686-5

	\vspace{9cm}

	\includegraphics[width=3cm,height=2.5cm,keepaspectratio]{Figures/nikhef-logo.eps}
	\hspace{2cm}
	\includegraphics[width=4.5cm,height=5cm,keepaspectratio]{Figures/nwo-logo.eps}
\end{center}

	\noindent This work is part of the research programme of the Foundation for Fundamental
	Research on Matter (FOM), which is part of the Netherlands Organisation for
	Scientific Research (NWO). It was carried out at the National Institute for
	Subatomic Physics (Nikhef) in Amsterdam.

\end{titlepage}
